%----------------------------------------
\begin{frame}
  \frametitle{Part 4}

{\Large\bf Worked examples}

\vspace{0.5cm}

{\large 1. Mutual exclusion --- a shared instrument}

\end{frame}

%----------------------------------------
\begin{frame}
  \frametitle{The system: a shared instrument} \small

\begin{block}{Informal requirements}\small
Two science instruments share a single power bus. Each must acquire the bus
before drawing power and release it afterwards.

\begin{enumerate}
\item {\bf Safety:} the two instruments are never powered simultaneously.
\item {\bf Liveness:} an instrument requesting power eventually gets it.
\item {\bf No deadlock:} the system never wedges.
\end{enumerate}
\end{block}

\pause
\vspace{0.3cm}

This is mutual exclusion. It is also, structurally, the Pathfinder bug,
the RTOS mutex, and half the concurrency in any flight system.

\end{frame}

%----------------------------------------
\begin{frame}[fragile]
  \frametitle{Attempt 1: a naive flag protocol} \small

\begin{lstlisting}[language=promela]
bool want[2];       /* want[i]: instrument i requests the bus */
byte powered = 0;   /* number of instruments currently powered */

active [2] proctype Instrument() {
  byte me = _pid;
  do
  :: want[me] = true;
     (!want[1-me]);          /* wait until the other is not asking */
     powered++;              /* CRITICAL: drawing power */
     assert(powered == 1);
     powered--;
     want[me] = false
  od
}
\end{lstlisting}

\pause

\begin{itemize}
\item \ipromela{_pid} is the process instance number, $0$ and $1$ here.
\item \ipromela{(!want[1-me])} is a {\bf blocking wait}, not an \ipromela{if}.
\item \ipromela{assert} states the safety requirement inline.
\end{itemize}

\end{frame}

%----------------------------------------
\begin{frame}[fragile]
  \frametitle{What \spin{} says} \small

\begin{lstlisting}[style=shell]
$ spin -a instrument.pml
$ gcc -o pan pan.c
$ ./pan

pan:1: invalid end state (at depth 4)
pan: wrote instrument.pml.trail

State-vector 20 byte, depth reached 5, errors: 1
\end{lstlisting}

\pause
\vspace{0.3cm}

\begin{block}{The counterexample --- a deadlock}\small
\begin{enumerate}
\item Instrument 0 sets \ipromela{want[0] = true}.
\item Instrument 1 sets \ipromela{want[1] = true}.
\item Instrument 0 blocks waiting for \ipromela{!want[1]}.
\item Instrument 1 blocks waiting for \ipromela{!want[0]}.
\end{enumerate}
Both are stuck. \spin{} found it at {\bf depth 4}, in milliseconds.
\end{block}

\pause

Note: the \ipromela{assert} never fired. Safety held; {\em progress} did not.

\end{frame}

%----------------------------------------
\begin{frame}[fragile]
  \frametitle{Attempt 2: Peterson's algorithm} \small

\begin{lstlisting}[language=promela]
bool want[2];
byte turn;
byte powered = 0;

active [2] proctype Instrument() {
  byte me = _pid;
  do
  :: want[me] = true;
     turn = 1 - me;                       /* politely defer */
     (!want[1-me] || turn == me);         /* wait */
     powered++;
     assert(powered == 1);                /* mutual exclusion */
     powered--;
     want[me] = false
  od
}
\end{lstlisting}

\pause
\vspace{0.2cm}

\begin{lstlisting}[style=shell]
$ ./pan
State-vector 20 byte, depth reached 27, errors: 0
\end{lstlisting}

{\bf Verified:} no assertion violation and no deadlock, over {\em all} interleavings.

\end{frame}
